\section{Curvature Estimation and Local Loss Change}
\label{appendix:loss-curvature}

Write $H(\theta)=\nabla^2\mathcal L(\theta)$ and assume the objective
is twice continuously differentiable along the plain-SGD step.
Integrate its gradient derivative on
$\theta_{t-2}-s\eta g_{t-1}$, add the endpoint-noise difference, and
project onto $-g_{t-1}/(\eta\|g_{t-1}\|^2)$. The result is
\begin{equation}
 \widehat\kappa_t=
 \int_0^1\frac{g_{t-1}^{\top}
 H(\theta_{t-2}-s\eta g_{t-1})g_{t-1}}
 {\|g_{t-1}\|^2}\,ds-\frac{\varepsilon_t}{\eta}.
 \label{eq:path-curvature}
\end{equation}
If $H$ is $L_H$-Lipschitz in operator norm on the step, its
full-batch specialization satisfies
\begin{equation}
 |\widehat\kappa_t-\kappa_t|
 \leq\int_0^1 L_Hs\eta\|g_{t-1}\|\,ds
 =\frac{L_H\eta}{2}\|g_{t-1}\|.
 \label{eq:curvature-error-bound}
\end{equation}
This bounds the difference between the measured path average and the
starting-point directional curvature. For Adam, integrate along its
adaptive step and project onto the preceding raw gradient. This gives
the exact path-average form of Eq.~\eqref{eq:adam-projected-change}.

Consider $\theta_{t-1}=\theta_{t-2}-\eta g_{t-1}$ and write
$\Delta\mathcal L_{t-1}=\mathcal L(\theta_{t-1})-\mathcal L(\theta_{t-2})$.
Integrating the full gradient along the step gives
\begin{equation}
\begin{aligned}
\Delta\mathcal L_{t-1}
=-\eta\int_0^1 g_{t-1}^\top
\nabla\mathcal{L}(\theta_{t-2}-s\eta g_{t-1})\,ds.
\end{aligned}
\label{eq:appendix-loss-integral}
\end{equation}
The trapezoidal approximation replaces the integrand by its endpoint
average. Substituting
$\nabla\mathcal L(\theta_{t-2})=g_{t-1}-\xi_{t-1}$ and
$\nabla\mathcal L(\theta_{t-1})=g_t-\xi_t$ yields
\begin{equation}
\begin{aligned}
\Delta\mathcal L_{t-1}
&\approx-\frac{\eta}{2}
\left(\|g_{t-1}\|^2+g_{t-1}^\top g_t\right)\\
&\quad+\frac{\eta}{2}
g_{t-1}^\top(\xi_{t-1}+\xi_t).
\end{aligned}
\label{eq:appendix-loss-gradient}
\end{equation}
The substitution $g_{t-1}^\top g_t=\|g_{t-1}\|\|g_t\|\rho_t$
recovers Eq.~\eqref{eq:loss-curvature}. The expression is exact on
quadratic objectives. A locally Lipschitz Hessian bounds its error
by a constant times $\|\eta g_{t-1}\|^3$.

Using $\rho_t\|g_t\|/\|g_{t-1}\|
\approx1-\eta\kappa_t+\varepsilon_t$ and
Eq.~\eqref{eq:epsilon} gives the equivalent curvature form,
\begin{equation}
\begin{aligned}
\Delta\mathcal L_{t-1}
&\approx-\eta\|g_{t-1}\|^2
\left(1-\frac{\eta\kappa_t}{2}\right)
\\&\quad+\eta\xi_{t-1}^\top g_{t-1}.
\end{aligned}
\label{eq:appendix-loss-curvature}
\end{equation}
It also follows directly from a second-order Taylor expansion at
$\theta_{t-2}$. With zero noise, the conditions for a loss-increasing
step reduce to Eqs.~\eqref{eq:loss-rise-alignment}
and~\eqref{eq:loss-rise-curvature}.

The noise term has a simple conditional interpretation. Given
$\theta_{t-2}$, unbiased sampling gives
$\mathbb E[\|g_{t-1}\|^2]=\|\nabla\mathcal L(\theta_{t-2})\|^2
+\mathbb E[\|\xi_{t-1}\|^2]$ and
$\mathbb E[\xi_{t-1}^\top g_{t-1}]=\mathbb E[\|\xi_{t-1}\|^2]$.
Their first-order variance terms cancel, leaving
\begin{equation}
\begin{aligned}
\mathbb E[\Delta\mathcal L_{t-1}\mid\theta_{t-2}]
&\approx-\eta\|\nabla\mathcal L(\theta_{t-2})\|^2\\
&\quad+\frac{\eta^2}{2}\mathbb E[
g_{t-1}^\top H_{t-1}g_{t-1}\mid\theta_{t-2}].
\end{aligned}
\label{eq:expected-loss-change}
\end{equation}
Thus sampling contributes through curvature at second order. A nonlinear
objective also contributes the finite-step remainder.
